% ============================================================
%  Extended Abstract — CEOBench — AI Agents Workshop
%  One page (letter), two columns, 10pt
% ============================================================
\documentclass[10pt,letterpaper,twocolumn]{article}

% ── Packages ────────────────────────────────────────────────
\usepackage[
  top=0.85in, bottom=0.85in,
  left=0.75in, right=0.75in,
  columnsep=0.25in
]{geometry}

\usepackage{times}           % Times New Roman font
\usepackage{microtype}       % Better text justification
\usepackage{titlesec}        % Section heading control
\usepackage{booktabs}        % Professional tables
\usepackage{graphicx}        % Figures
\usepackage{amsmath,amssymb} % Math
\usepackage{hyperref}        % Clickable refs & DOI links
\usepackage{xcolor}          % Colors for hyperlinks
\usepackage{enumitem}        % Compact lists
\usepackage{cite}            % Compressed numeric citations
\usepackage{fancyhdr}        % Header / footer

% ── Hyperlink colours ───────────────────────────────────────
\hypersetup{
  colorlinks=true,
  linkcolor=black,
  citecolor=black,
  urlcolor=blue
}

% ── Section heading sizes ───────────────────────────────────
\titleformat{\section}{\normalsize\bfseries\uppercase}{\thesection.}{0.5em}{}
\titleformat{\subsection}{\normalsize\bfseries\itshape}{\thesubsection.}{0.5em}{}
\titlespacing*{\section}{0pt}{5pt}{1pt}
\titlespacing*{\subsection}{0pt}{3pt}{1pt}

% ── Paragraph & list spacing ────────────────────────────────
\setlength{\parskip}{1.5pt}
\setlength{\parindent}{1em}
\setlist{noitemsep, topsep=1pt, leftmargin=1.2em}

% ── Header / Footer ─────────────────────────────────────────
\pagestyle{fancy}
\fancyhf{}
\renewcommand{\headrulewidth}{0.4pt}
\fancyhead[L]{\small\textit{NE Agents Day 2026}}
\fancyhead[R]{\small Extended Abstract}
\fancyfoot[C]{\small\thepage}

% ============================================================
\begin{document}
% ============================================================

% ── Title block (spans both columns) ────────────────────────
\twocolumn[{%
  \begin{center}
    {\Large\bfseries \textsc{CEOBench}: Evaluating Language Model Agents\\[2pt] as Long-Horizon Strategic Decision-Makers\par}
    \vspace{5pt}
    {\normalsize
      \textbf{Haozhe Chen}\textsuperscript{1}\quad
      \textbf{Zhuang Liu}\textsuperscript{1}
    \par}
    \vspace{3pt}
    {\small
      \textsuperscript{1}Department of Computer Science, Princeton University
    \par}
    \vspace{3pt}
    {\small
      \texttt{\{hc5019, zhuangl\}@princeton.edu}
    \par}
  \end{center}

  \vspace{3pt}
  \noindent\textbf{Keywords:} AI agents, long-horizon planning, partial observability, non-stationary environments, strategic decision-making
  \vspace{8pt}
}]

% ============================================================
%  BODY
% ============================================================

\section{Introduction and Related Work}

Language model agents have demonstrated proficiency at isolated, short-horizon tasks---resolving GitHub issues~\cite{jimenez2024swebench}, handling customer service requests~\cite{yao2024taubench}, and navigating websites~\cite{zhou2024webarena}. As agents approach reliable execution of such individual tasks, a natural next question is what we should expect them to do after the local task is no longer the bottleneck. A crucial part of what we consider great human intelligence---the kind behind many major human achievements---is the ability to navigate uncertainty over long horizons, read between the lines of incomplete information, adapt when the world shifts, and orchestrate many moving parts toward a coherent goal. This remains largely untested.

Existing benchmarks evaluate episodic tasks in stationary environments where the horizon spans minutes to hours~\cite{jimenez2024swebench,zhou2024webarena,liu2023agentbench,mialon2024gaia}. Even $\tau$-bench~\cite{yao2024taubench}, which models customer service as a POMDP, is scoped to single conversations with no persistent consequences. VendingBench~\cite{backlund2025vendingbench} takes a step toward long-horizon evaluation, but the action space is narrow and the environment lacks the interdependent decision domains and latent-state complexity that require coordinated strategic planning. We introduce \textsc{CEOBench} to fill this gap.

\section{Benchmark Design}

In \textsc{CEOBench}, an agent operates a subscription-based software company over 500 simulated days, taking hundreds of actions per day across six interdependent domains: pricing, infrastructure capacity, marketing, R\&D, enterprise negotiations, and financial management. It accesses 34 tools and 19 database tables. The benchmark stresses four axes of difficulty:
\begin{itemize}
  \item \textbf{Long-horizon planning.} A 500-day episode spanning tens of thousands of actions requires a coherent long-term strategy balancing short-term survival against long-term growth.
  \item \textbf{Partial observability.} Customer satisfaction, reputation, churn probabilities, and preference drift parameters are hidden---the agent infers latent state from aggregate metrics and social media.
  \item \textbf{Non-stationarity.} Preferences drift monthly, macroeconomic conditions follow a stochastic mean-reverting process, and competitor events trigger exogenous demand shocks.
  \item \textbf{Multi-capability coordination.} Decisions are tightly coupled: a capacity shortfall degrades quality, eroding reputation, reducing acquisition, lowering revenue, and constraining future investment.
\end{itemize}

\vspace{-3pt}
Feedback is sparse: revenue arrives on 30-day billing cycles, R\&D takes 5--30 days, and reputational damage may surface weeks later. The score is final cash balance, with bankruptcy ending the episode early.

\section{Preliminary Results}

We evaluate frontier models including Claude Sonnet 4.6 and GLM-5 over 500-day episodes (full benchmark duration). Preliminary runs reveal wide performance variance across seeds and models, with agents frequently failing to adapt to mid-episode distributional shifts. Common failure modes include over-investing in R\&D without securing short-term revenue, failing to detect preference drift until churn cascades, and mismanaging capacity during growth. Detailed analysis is forthcoming.

\section{Conclusion}

\textsc{CEOBench} provides a challenging testbed for evaluating whether language model agents can move beyond competent task execution toward sustained, adaptive, strategic reasoning---the kind of intelligence we recognize in people but have not yet measured in machines.

% ── References ──────────────────────────────────────────────
{\small
\bibliographystyle{abbrv}
\bibliography{references}
}

\end{document}
